\lesson{4}{Nov 8 2021 Mon (13:20:50)}{Influences Public Policy}{Unit 3}

\subsubsection*{Influencing Public Policy}

Every person and organization in the \bf{United States} has the power to
influence what public policies the government enacts and enforces.
\bf{Representatives} at the:

\begin{itemize}
    \item Local
    \item State
    \item National Levels
\end{itemize}

seek to represent the wishes of as many of their constituents as possible, while
taking care not to restrict the rights of those with less popular views.
Therefore, it is true that there is strength in numbers in a \bf{Representative
Democracy}.

People tend to seek out others who have had similar experiences or have similar
views about the world. They tend to share a common purpose, and that purpose
could be a change in the laws. Groups that tend to have a strong influence on
government in the \bf{United States} include:

\begin{itemize}
    \item Political Parties
    \item Special Interest Groups
    \item Media
\end{itemize}

\subsubsection*{One Individual Having Power}

Americans have power as individuals in many different ways. Voting in elections
is probably the most obvious method. People can choose to run for \bf{Elected
Office} themselves, or they can service staff members and government offices
responsible for communicating with the public or drafting new public policies.
Individuals can share their opinions through the media or contacting public
officials. They can speak at meetings of lawmaking bodies or participate in
surveys that help representatives learn about constituent needs. History has
shown that one person's experience can have a profound effect on public policy.

An example is the \bf{AMBER alert system}, a nationwide voluntary program where
law enforcement agencies at all levels of government work together to locate a
missing child as quickly as possible. \bf{Congress} initiated the program as a
legacy to nine-year-old \bf{Amber Hagerman}, a Texas girl who, in $1996$, was
kidnapped while riding her bike and murdered. Stories like \bf{Amber's} can have
a strong emotional impact, persuading lawmakers to make new policies and
programs in the hope that they can prevent similar tragedies in the future.

The individual power of one person can change over time. For example, \bf{State
Legislatures} originally chose the individuals who would serve as senators in
the \bf{U.S. Congress}. The \it{Seventeenth Amendment} was added to the
\bf{Constitution} in $1913$, changing the election procedure to allow American
voters to choose their senators by direct election. The individual voting power
of the average American was increased through this amendment.

\subsubsection*{Ideology Affecting Political Parties}

Individuals also choose to influence government by joining their voices with
like-minded people. For example, a person may join a political party that shares
his or her beliefs about the purpose of government and how it should function.
This set of beliefs is referred to as a political ideology. Political parties
share a common political ideology and have great influence on who is elected to
political office and the policies that lawmakers create. Major parties have
chapters that coordinate their efforts at the local and state, as well as
national, levels.

The United States has two major political parties—the Democratic Party and the
Republican Party. Each party has a platform that describes its beliefs, values,
and ideas. Members of a political party usually agree with the party platform,
though they might have a different view on certain issues.

Americans are most familiar with the two main parties and usually identify
themselves as a Democrat or Republican. However, there are other political
parties. We refer to them as third parties. Third parties often form because a
group of people believe that the main parties fail to address or deal with a
certain issue. That issue or problem becomes the focus of the third party,
leading to the related term "single issue" parties for that reason. An example
is the Green Party of the United States whose focus is protecting and preserving
the natural environment. Its members believe that the way big businesses grow
and operate is harming Earth and U.S. society in the long term. Some people, of
any ideology, register to vote or run for political office as an independent.
That means they did not register as a member of a political party.

\paragraph{Political Spectrum}

\begin{itemize}
    \item "To the left," "Center," "To the Right"   \\
        \bf{Ideology is very complex. Dividing into the labels of liberal, moderate, and conservative are helpful to describe political ideology, though sometimes it can be misleading because there is much room for variation. For example, two people might be "conservative," but one could be a bit more conservative than the other. Have you ever heard a person's views described as "to the left" or "right of center"? "To the left" refers to liberal views, while "to the right" refers to conservative views. We sometimes describe moderates as "in the center" or "on the fence" because they could lean in either direction. You can also label political parties by their beliefs.}
    \item Liberal                                   \\
        \bf{Liberals generally believe that it is the government's responsibility to ensure the people's well-being. They tend to increase government programs and taxes to pay for them. For example, increasing funding for environment programs or assistance for families in poverty tend to be liberal ideas. In the United States, voters who would describe themselves as liberal tend to vote for candidates that are members of the Democratic Party.}
    \item Moderate                                  \\
        \bf{Moderates tend to be liberal on some issues and conservative on others. Most U.S. voters are moderate, or they are mostly moderate but might tend to lean to one side more often. Candidates for political office compete hardest for moderate votes. Candidates for office need to capture the votes of as many moderate people as possible. Because most people are moderate, a political party platform will stay away from extreme views if they want to increase their support. Whom a moderate will vote for depends on the specific issues in that election and the candidates' plans to address those issues.}
    \item Conservative                              \\
        \bf{Conservatives generally believe that it is the people's own responsibility to ensure their well-being. They tend to limit government spending on programs and raising taxes. For example, prioritizing government funds on national security and the military tend to be conservative ideas. In the United States, voters who would describe themselves as conservatives tend to vote for candidates that are members of the Republican Party.}
    \item Democratic Party                          \\
        \bf{The Democratic Party is described as "liberal," but in reality members are typically somewhere between liberal and moderate on this spectrum. The major political parties avoid extreme views in ideology so they can capture as much support as possible from voters.}
    \item Republican Party                          \\
        \bf{The Republican Party is described as "conservative," but in reality, members are typically somewhere between conservative and moderate on this spectrum. The major political parties avoid extreme views in ideology so they can capture as much support as possible from voters.}
\end{itemize}

\subsubsection*{Political Parties Developing in the United States}

The name of a political party is just that–a name. To understand the party you
have to focus on its platform, which can change over time. The \bf{Constitution}
does not mention political parties, but they developed out of competing
interests during the \bf{Founding Era}. Political parties have existed ever
since, transitioning in response to new challenges and the emergence of powerful
leaders.

Even the earliest political parties in the \bf{United States} took positions on
a wide variety of issues. Their broad focus and avoidance of extremism left
little room for any third party to come along and distinguish itself. The
principle of majority rule in American elections lends itself to the two-party
political system that developed. You can either support a policy or issue, or
choose not to support it.

Political historians generally agree on five major party systems in American
history, meaning that each new system reveals shifts in the platform and
ideology of each party and their bases of voter support. The name they gave to
each shift is "realignment."

\paragraph{George Washington Era and Rise of the First Party System} When George
Washington was \bf{President}, Thomas Jefferson and Alexander Hamilton were
members of his cabinet and often at odds with one another. Jefferson was an
Anti-Federalist and Hamilton a Federalist, from each of the two groups that
arose during the Constitution's ratification process. Jefferson, frustrated that
Washington generally took Alex Hamilton's side, resigned his position. Concerned
about their feuding and the example it could set for the future of American
politics, Washington warned the nation against the creation of "factions" in his
farewell address. He thought that political parties could undermine the purpose
of the Republic he helped create–to represent the will of all people, not just
those with certain views.

\paragraph{1796 John Adams} Many important national leaders agreed with
Jefferson's views, however, so he was not out of the public spotlight for long.
The year 1796 included the first Presidential election, with a clear distinction
of candidates and their followers. Jefferson supporters, many who had been
Anti-Federalists, now called themselves Democratic-Republicans. Federalists,
retaining their ratification-era name, supported John Adams. Adams won the
election and Jefferson became the vice \bf{President}, because the original
Electoral College system gave the position to the candidate with the second-most
votes. Federalist votes for a second candidate were scattered among several
people, leaving Jefferson with the next highest number of electoral votes.

\paragraph{1800 Thomas Jefferson} Determined to prevent a "divided government"
in the Election of 1800, Jefferson and Adams chose "running mates," encouraging
their followers to vote for a specific person who would become vice
\bf{President} from the same party. Jefferson's following had grown and he
defeated the incumbent Adams in 1800. This election marks the first time in
American history that the presidency transferred from one political ideology to
another.

\paragraph{1824 Andrew Jackson and Rise of the Second Party System} The
Democratic-Republicans continue to dominate the presidency until 1824, when they
split into two separate parties. Political participation by the public had grown
exponentially, fueled by the change in many states to allow voters to select who
would run as candidate for \bf{President}. Before, the state legislatures chose
the candidate. Several major candidates ran in this election, including the war
hero Andrew Jackson. He won the popular vote largely based on gathering support
from people in the South and West who wanted to expand the country's territory
and remove Native Americans. The economy and expansion were major issues of the
time. However, Jackson lost in the Electoral College. His followers were called
the Jacksonian-Democrats, later simply Democrats, who did vote him into the
presidency in 1828. The party that emerged as Jackson's opposition was the Whig
Party, whose members wanted to encourage economic growth through national
policies that Jackson considered unconstitutional.

\paragraph{1860 Abraham Lincoln and Rise of the Third Party System} By the
mid-1800s, the dominant issue in American politics was slavery. A third party
emerged that was against slavery, called the Republican Party. Support for this
party was concentrated in states that had abolished slavery, especially in the
North. The original Jacksonian-Democrats had split into separate parties since
the issues that had once unified them were now obsolete or minor in comparison
to the current politics. Abraham Lincoln won the Electoral College in 1860 as
the Republican candidate. No Southern states gave electoral votes to Lincoln,
but their Democratic votes were split among several candidates. The Republican
Party of Lincoln is the root of the modern Republican Party, and they supported
many of the economic policies the Whigs had once encouraged, such as railroad
building and high tariffs, which are taxes on goods imported from other
countries–liberal policies. However, their power in making policy was tempered
by the fact that there were still many conservative Democrats in Congress. Note
that the ideology of the Republicans at that time was quite different from
modern Republicans.

\paragraph{1896 William Jennings Bryan and Rise of the Fourth Party System} The
Democratic candidate in 1896, William Jennings Bryan, alienated Northern
Democrats by supporting policies that would mostly benefit the interests of
rural farmers concentrated in the South and West. Instead of slavery dividing
the major parties, now the big issue was economic policies like currency and
tariffs. Many Northern Democrats shifted their votes to Republican candidate
William McKinley, who pushed for policies that would benefit cities and
industry. The party names remained the same, but their membership and leadership
shifted along with ideology and platform.

\paragraph{1932 Franklin D. Roosevelt and Rise of the Fifth Party System} In
1932, the United States was dealing with the Great Depression, a severe downturn
in the economy. Once again, the economy would cause a shift in the parties.
Franklin D. Roosevelt, as the Democratic candidate, promised Americans a "New
Deal" that would bring relief to struggling Americans from many different
backgrounds, from cities and farms alike. He won the election and, as
\bf{President}, instituted many federal programs that would provide assistance
and put people to work. The costs were high, however, and those who thought new
deal policies gave the federal government too much power rallied behind the
Republican Party, adopting views that are more conservative.

\paragraph{2008 Barack Obama and … What?} In the second half of the 20th
century, a major change to note is that the South tended to vote more Republican
than other regions of the nation. Other than that, analysts describe a period of
decay for the political party. More voters than ever are registering as
independents rather than as members of a major party. The two major parties
sometimes have views that seem indistinguishable from one another in an attempt
to capture the moderate voter. Third parties have had notable impact in several
state and federal elections. Several Presidents have faced congresses dominated
by the opposite party. In 2008, Democrat Barack Obama won the presidency,
carrying many states that traditionally voted Republican. Analysts debate
whether this indicates a new party realignment or simply further decay. It's too
early to tell.

\subsubsection*{Political Parties Affecting Public Policy}

The main goal of a political party is to run campaigns to get its candidates
elected so they can affect public policy. The political party with the majority
of seats in \bf{Congress} has greater control over the direction of public
policy that passes into law. The majority leaders in both the \bf{House} and the
\bf{Senate} have some power over what bill ideas go to the committees for
consideration. Lawmaking happens through majority votes of the legislators, so
policies favored by the majority party are more likely to survive the process.

If the majority party is the same as the \bf{President's} party, the
\bf{President} will have an easier time persuading \bf{Congress} to follow their
policy agenda. \bf{Presidents} are considered the leader of the political party
to which they belong. Presidents who have served during the time when the
opposite party held the majority in Congress have found it more difficult to
gather the support they needed. During such a time, it is more common for a
\bf{President} to use the veto power.

\bf{State Governments} show similar trends between governors and state
legislatures. State-level policies can provide a strong argument for or against
a particular federal policy idea, depending on whether the program or law is
successful in that state and which party supported it.

\subsubsection*{Special Interest Groups}

Like political parties, people form special interest groups because they share a
common idea or experience and want to influence public policy. However, interest
groups focus on a smaller range of issues, which usually relate to one another,
and try to persuade those in office to enact policies that favor their position
on those issues. Typically, a person will join a special interest group because
they feel very passionate about the issue or issues addressed by the group.

One example of a special interest group is a labor union, which focuses on the
concerns of workers, often in a particular field such as education. A prominent
environmental interest group is the Sierra Club, which takes positions on many
important national issues related to resource conservation and climate. A
business corporation or an institution, such as a police department or school
district, can even operate as a special interest group when a legislative body
is considering a policy that would affect it.

Interest groups seek to inform the American public, as well as persuade
lawmakers and other officials. The more support an interest group has, the more
likely it will be able to convince lawmakers to pass or not pass a particular
policy idea. Interest groups can form, dissolve, and change their policy focus
depending on what issues are at the center of attention at any point in time.

Interest groups might hire lobbyists to meet with lawmakers. Many Americans
think of lobbying as a negative term. They feel that special interests have too
much influence over lawmakers and do not represent the majority of the public.
To address these concerns, today's lobbyists must follow strict ethics rules.
Examples include several rules on when giving a gift to a government official is
acceptable.

\newpage
